\section{Revisão Bibliográfica}
\label{sec:revisao}

\subsection{Arquiteturas de recarga de veículos elétricos}
Os carregadores de veículos elétricos podem ser classificados em embarcados (\textit{on-board},OBC) que operam tipicamente entre 3,7 kW e 22 kW e a externos (\textit{off-board})
de dezenas a centenas de kW \cite{yilmaz2013review,tu2019extreme}. Ambos tipos 
apresentam dois estágios, o primeiro é um retificador PFC que é responsável por regular 
o barramento CC intermediário e o seguindo é um conversor CC-CC isolado que adapta a tensão do barramento CC à tensão da bateria (tipicamente entre 400 e 800 V). 

Estações de recarga rápidas apresentam o estágio CA-CC formado por módulos de 15 a 30 kW em paralelo para diminuir o esforço das chaves dos retificadores, sendo a escolha da 
topologia utilizada dependente da eficiência, densidade de potência e qualidade de energia necessária. Como os 
OBCs podem operar tanto em redes monofásicas quanto em redes trifásicas, topologias com
o ponto médio do barramento ligado aso neutro da rede são comumente utilizadas \cite{schafmeister2021transformerless}.
\subsection{Topologias de retificadores trifásiocos}

Existem diversas topologias de retificadores trifásicos, entre elas podem ser destacadas algumas características. 

As pontes de diodos com conversor boost são simples e necessitam de apenas uma chave, por conta dessa particularidade, apresentam
apenas um nível de tensão e um fluxo de corrente unidirecional, não sendo possível controlar o FP, o que provoca
uma corrente de linha distorcida. As topologias de retificador com PWM de dois níveis apresentam seis chaves e são bidirecionais, além de apresentarem
dois níveis de tensão, permitindo o controle completo do fator de potência e da corrente de linha, porém apresentam uma maior complexidade de controle e suas chaves bloqueiam a tensão do barramento, assim ficam sobre um estresse maior \cite{kolar2013essence}. 

Já as topologias NPC/tipo-T apresentam doze chaves e 3 níveis de tensão, além de permitirem um fluxo de corrente
bidirecional, o que torna possível a operação em modo de recarga e de injeção de energia na rede. Apesar de reduzirem os esforçoes sobre cada chave, seu custo é maior
tendo em vista a necessidade de doze chaves \cite{safayatullah2022comprehensive}. O retificador de Vienna de três fios \cite{kolar1997novel} é a versão unidirecional do 
tipo-T, ele apresenta três chaves bidirecionais e seis diodos rápidos que produzem três níveis de tensão, com bloqueio de tensão em $V_{\sub{dc}}/2$. Esse comportamento garante uma alta eficiência e alta densidade de potência, porém não permite a operação bidirecional 
e o controle de fator de potência é limitado. O retificador de Vienna de 4 fios foi proposto, especialmente no contexto de carregamento de veículos elétricos, tendo em vista sua aplicação em estações de recarga rápidas fora do veículo \cite{senthilkumar2025iipq} e
para carregadores embarcados sem transformados \cite{schafmeister2021transformerless}. Assim como o retificador de Vienna de três fios, o retificador de Vienna de quatro fios apresenta três níveis de tensão, porém com a adição de um quarto fio, o neutro da rede, é possível controlar o fator de potência da rede, seu fluxo de corrente é unidirecional e apresenta 3 níveis de tensão.

\subsection{Estratégias de controle}

Para um retificador de Vienna de quatro fios, os mecanismos doe controle, de forma geral, se concentram em três malhas, uma para a corrente de entrada,
uma para a tensão do barramento CC e uma para o equilíbrio do ponto médio do barramento CC. A malha de corrente de entrada é responsável por controlar a corrente de linha, garantindo que ela siga uma referência senoidal em fase com a tensão da rede, o que permite o controle do fator de potência. A malha de tensão do barramento CC regula a tensão do barramento intermediário, garantindo que ela permaneça dentro dos limites especificados para a operação segura do conversor. Por fim, a malha de equilíbrio do ponto médio do barramento CC assegura que as tensões nos capacitores do barramento permaneçam equilibradas, evitando sobrecargas em um dos capacitores e garantindo a estabilidade do sistema.

Para o retificador destacado anteriormente, o equilíbrio é realizado com um termo aditivo na amplitude da referência, com sinal dependente do semiciclo \cite{alahuhtala2006fourwire}.

